\section[Cote 19 : le théorème de comonadicité 1.12, la comonade sur une base produit et le théorème de Gabriel--Popescu]{Cote 19 : le théorème de comonadicité 1.12, la comonade sur une base produit et le théorème de Gabriel--Popescu\protect\verifie{Folder19}}\label{sec:19}

La cote 19, \emph{Topos : notes manuscrites (s.d.)}, quatre-vingt-quinze pages que l'inventaire date \q{[à partir de 1958-1973]} et range dans le groupe Théorie des motifs, est un dossier de travail en neuf liasses. La première, « Tapis cartésiens » (pages~3 à~11), étudie la comonade d'une adjonction et ses coalgèbres ; la page~4 énonce le \q{Théorème 1.12}, suivi d'un nom propre que la transcription lit \q{Beck} avec doute, écrit \q{il faut et il suffit} pour~b), et porte en marge la restriction \q{il suffit de regarder les doubles flèches $u, v$ de $A$ telles que $\mathrm{Ker}(f(u), f(v))$ existe \quad (condition C)} (\lot{19}{1}{4}). La deuxième, « Cohomologie relative », s'ouvre pages~14 et~15 sur la représentabilité : pour une catégorie à générateur $U$, elle affirme une équivalence avec les modules sur $\mathrm{End}(U)$, \q{[détails laissés au lecteur]} (\lot{19}{1}{15}).

\noindent\emph{Conventions.} On note, comme la lecture, $f : A \to B$ l'adjoint à gauche de $g : B \to A$, $\eta : \mathrm{id}_A \to gf$ l'unité, $\varepsilon : fg \to \mathrm{id}_B$ la coünité, et $\varphi = fg$ la comonade, de comultiplication $\delta_X = f(\eta_{gX})$. Une $\varphi$-coalgèbre est un couple $(X, a)$, $a : X \to \varphi X$, tel que $\varepsilon_X \circ a = 1_X$ et $\delta_X \circ a = \varphi(a) \circ a$ ; un morphisme $(X, a) \to (X', a')$ est une flèche $t : X \to X'$ telle que $a' \circ t = \varphi(t) \circ a$. Le foncteur de comparaison $h : A \to \varphi\text{-}\mathbf{Coalg}$ envoie $Y$ sur $(fY, f(\eta_Y))$ et $\psi$ sur $f(\psi)$. Le \emph{noyau} $\mathrm{Ker}(u, v)$ d'une double flèche est son égalisateur, et $f$ \emph{y commute} s'il envoie un noyau de $(u, v)$ sur un noyau de $(fu, fv)$. On utilise les identités triangulaires $\varepsilon_{fY} \circ f(\eta_Y) = 1$ et $g(\varepsilon_X) \circ \eta_{gX} = 1$, et la bijection d'adjonction $t \mapsto t^\flat = g(t) \circ \eta_Y$, de $\operatorname{Hom}_B(fY, X)$ sur $\operatorname{Hom}_A(Y, gX)$, d'inverse $s \mapsto \varepsilon_X \circ f(s)$.

\begin{definition}
Le foncteur $f$ vérifie la \emph{condition~C} si, pour toute double flèche $(u, v)$ de $A$ telle que $\mathrm{Ker}(fu, fv)$ existe, $\mathrm{Ker}(u, v)$ existe et $f$ y commute. Un triplet $(e, s, t)$ formé de $e : E \to Y$, $s : Y \to E$ et $t : Z \to Y$ est un \emph{noyau scindé} de $u, v : Y \rightrightarrows Z$ si $ue = ve$, $se = 1_E$, $tu = 1_Y$ et $tv = es$. Une double flèche $(u, v)$ de $A$ est \emph{$f$-scindée} si $(fu, fv)$ admet un noyau scindé.
\end{definition}

\begin{enonce}[Théorème 1.12, page 4]
Soit $f \dashv g$ une adjonction, $\varphi = fg$, et $h$ le foncteur de comparaison.
\begin{enumerate}
\item[a)] Si les noyaux de doubles flèches existent dans $A$ et si $f$ y commute, alors $h$ est essentiellement surjectif.
\item[b)] $h$ est une équivalence si et seulement si $f$ est conservatif et si, pour toute double flèche $(u, v)$ de $A$ telle que $\mathrm{Ker}(fu, fv)$ existe, $\mathrm{Ker}(u, v)$ existe et $f$ y commute.
\end{enumerate}
La restriction portée en marge --- il suffit de considérer les doubles flèches $(u, v)$ dont l'image par $f$ a un noyau --- est la condition~C.
\end{enonce}

\begin{lemme}\label{lem:19-scinde}
Si $(e, s, t)$ est un noyau scindé de $(u, v)$, alors $e$ est un noyau de $(u, v)$, et l'image de $(e, s, t)$ par tout foncteur est un noyau scindé de l'image de $(u, v)$.
\end{lemme}

\begin{proof}
Soit $z : T \to Y$ avec $uz = vz$. Alors $z = tuz = tvz = e(sz)$, donc $z$ se factorise par $e$, et de façon unique, puisque $ez' = ez''$ entraîne $z' = sez' = sez'' = z''$. Les quatre conditions sont des égalités entre composés, qu'un foncteur conserve.
\end{proof}

\begin{lemme}\label{lem:19-paire}
Pour toute coalgèbre $(X, a)$, la double flèche $(\eta_{gX}, g(a)) : gX \rightrightarrows gfgX$ est $f$-scindée : $(a, \varepsilon_X, \varepsilon_{fgX})$ est un noyau scindé de $(f(\eta_{gX}), fg(a))$. Elle est de plus coréflexive : $g(\varepsilon_X)$ est une rétraction commune de $\eta_{gX}$ et de $g(a)$.
\end{lemme}

\begin{proof}
On a $f(\eta_{gX}) \circ a = \delta_X \circ a = \varphi(a) \circ a = fg(a) \circ a$ par coassociativité ; $\varepsilon_X \circ a = 1_X$ par l'axiome de coünité ; $\varepsilon_{fgX} \circ f(\eta_{gX}) = 1$ par l'identité triangulaire en $gX$ ; et $\varepsilon_{fgX} \circ fg(a) = a \circ \varepsilon_X$ par naturalité de $\varepsilon$. Enfin $g(\varepsilon_X) \circ g(a) = g(\varepsilon_X \circ a) = 1$ et $g(\varepsilon_X) \circ \eta_{gX} = 1$.
\end{proof}

\begin{lemme}\label{lem:19-comp}
Supposons que, pour toute coalgèbre $(X, a)$, la double flèche $(\eta_{gX}, g(a))$ ait un noyau $e_X : R(X, a) \to gX$ et que $f$ y commute.
\begin{enumerate}
\item Pour tout $Y \in A$, $t \mapsto t^\flat$ induit une bijection entre les morphismes de coalgèbres $hY \to (X, a)$ et les flèches $k : Y \to R(X, a)$, par $t^\flat = e_X \circ k$.
\item Pour toute coalgèbre $(X, a)$, la flèche $\theta = \varepsilon_X \circ f(e_X)$ est un isomorphisme de coalgèbres de $h(R(X, a))$ sur $(X, a)$ ; en particulier $h$ est essentiellement surjectif.
\item Si de plus $f$ est conservatif, $h$ est une équivalence.
\end{enumerate}
\end{lemme}

\begin{proof}
(1) Soit $t : fY \to X$. Les transposées des deux membres de $a \circ t = \varphi(t) \circ f(\eta_Y)$ sont $g(a) \circ t^\flat$ et, puisque $\varphi(t) \circ f(\eta_Y) = f(t^\flat)$, $gf(t^\flat) \circ \eta_Y = \eta_{gX} \circ t^\flat$ par naturalité de $\eta$. La transposition étant bijective, $t$ est un morphisme de coalgèbres si et seulement si $t^\flat$ égalise $(\eta_{gX}, g(a))$, c'est-à-dire se factorise, de façon unique, par $e_X$.

(2) Posons $E = R(X, a)$. Par~(1) appliqué à $k = 1_E$, la flèche $\theta$, de transposée $e_X$, est un morphisme de coalgèbres $hE \to (X, a)$, et $a \circ \theta = \varphi(\theta) \circ f(\eta_E) = f(\theta^\flat) = f(e_X)$. Or $a$ est un noyau de $(f(\eta_{gX}), fg(a))$ par les lemmes~\ref{lem:19-scinde} et~\ref{lem:19-paire}, et $f(e_X)$ en est un par hypothèse ; donc $\theta$ est un isomorphisme. Son inverse est un morphisme de coalgèbres :
\[
\varphi(\theta^{-1}) \circ a = \varphi(\theta^{-1}) \circ a \circ \theta \circ \theta^{-1} = \varphi(\theta^{-1}) \circ \varphi(\theta) \circ f(\eta_E) \circ \theta^{-1} = f(\eta_E) \circ \theta^{-1}.
\]

(3) Soit $Y' \in A$. La flèche $\eta_{Y'}$ égalise $(\eta_{gfY'}, gf(\eta_{Y'}))$, par naturalité de $\eta$ ; c'est la double flèche attachée à la coalgèbre $hY' = (fY', f(\eta_{Y'}))$, et l'on a donc $\eta_{Y'} = e_{hY'} \circ k_{Y'}$ pour une flèche $k_{Y'} : Y' \to R(hY')$. Appliquons $f$ : $f(\eta_{Y'})$ est un noyau de l'image de cette double flèche (lemmes~\ref{lem:19-scinde} et~\ref{lem:19-paire}, avec $a = f(\eta_{Y'})$), $f(e_{hY'})$ aussi par hypothèse, et $f(e_{hY'}) \circ f(k_{Y'}) = f(\eta_{Y'})$ ; donc $f(k_{Y'})$ est un isomorphisme, et $k_{Y'}$ aussi, $f$ étant conservatif. Pour $\psi : Y \to Y'$, la transposée de $h(\psi) = f(\psi)$ est $gf(\psi) \circ \eta_Y = \eta_{Y'} \circ \psi = e_{hY'} \circ (k_{Y'} \circ \psi)$ : dans la bijection~(1), $h(\psi)$ correspond à $k_{Y'} \circ \psi$. Comme $\psi \mapsto k_{Y'} \circ \psi$ est bijective, $h$ est pleinement fidèle ; il est essentiellement surjectif par~(2), c'est donc une équivalence.
\end{proof}

\begin{proof}[Démonstration de a) et de la suffisance de b)]
Si $f$ vérifie la condition~C, chaque double flèche $(\eta_{gX}, g(a))$ a un noyau auquel $f$ commute, puisque son image a un noyau (lemmes~\ref{lem:19-paire} et~\ref{lem:19-scinde}). Le lemme~\ref{lem:19-comp}\,(2) donne alors a) sous la forme de la marge, et la forme du texte s'en déduit, car si tous les noyaux existent dans $A$ et que $f$ y commute, la condition~C est vérifiée. Si de plus $f$ est conservatif, $h$ est une équivalence par le lemme~\ref{lem:19-comp}\,(3).
\end{proof}

\begin{proof}[La nécessité de b) est fausse]
Soit $B$ la catégorie des couples d'ensembles $(X, Y)$, et $A$ la sous-catégorie pleine des couples tels que $Y \neq \varnothing$ entraîne $X \neq \varnothing$ ; $f$ est l'inclusion. Posons $g(X, Y) = (X, Y)$ si $X \neq \varnothing$ et $g(X, Y) = (\varnothing, \varnothing)$ sinon ; $g(X, Y)$ est dans $A$. Pour $(X', Y') \in A$, une flèche $(p, q) : (X', Y') \to (X, Y)$ de $B$ se factorise de façon unique par $g(X, Y) \to (X, Y)$ : si $Y' \neq \varnothing$, alors $X' \neq \varnothing$, donc $X \neq \varnothing$ ; si $Y' = \varnothing$, $q$ est l'application vide. Ainsi $f \dashv g$, et $f$ est une inclusion coréflexive. Une telle inclusion est pleinement fidèle, donc conservative, et comonadique : c'est classique --- la comonade est idempotente, et ses coalgèbres sont, à isomorphisme près, les objets de $A$. Le foncteur $h$ est donc une équivalence.

Mais $f$ ne vérifie pas la condition~C. Soient $u, v : (\ast, \ast) \rightrightarrows (\{0, 1\}, \ast)$ les flèches de $A$ qui envoient le point de la première composante sur $0$ et sur $1$ et sont l'identité sur la seconde. Leur noyau dans $B$ existe, c'est $(\varnothing, \ast)$. Si $K \to (\ast, \ast)$ était un noyau de $(u, v)$ dans $A$ auquel $f$ commute, $K$ serait isomorphe à $(\varnothing, \ast)$, de seconde composante non vide ; étant dans $A$, il aurait un point $x$ dans sa première composante, et $u$ et $v$ l'enverraient sur le même élément, c'est-à-dire $0 = 1$ : c'est absurde. Le noyau de $(u, v)$ dans $A$ est $(\varnothing, \varnothing)$, et $f$ n'y commute pas.
\end{proof}

\begin{enonce}[Théorème de Beck, forme exacte]
Le foncteur $h$ est une équivalence si et seulement si $f$ est conservatif et si, pour toute double flèche $f$-scindée $(u, v)$ de $A$, $\mathrm{Ker}(u, v)$ existe et $f$ y commute.
\end{enonce}

\begin{proof}
La condition est suffisante par les lemmes~\ref{lem:19-paire} et~\ref{lem:19-comp}\,(3), qui n'utilisent que des doubles flèches $f$-scindées. Elle est nécessaire. Si $f(\psi)$ est inversible, $h(\psi)$ est un morphisme de coalgèbres dont la flèche sous-jacente est inversible, donc un isomorphisme de coalgèbres (le calcul de~(2) ci-dessus), et $\psi$ est inversible parce que $h$, pleinement fidèle, reflète les isomorphismes. Que le foncteur d'oubli des coalgèbres, donc $f$, crée les noyaux des doubles flèches $f$-scindées est la partie classique du théorème de Beck, que la preuve vérifiée invoque.
\end{proof}

\noindent\emph{Une base produit.} Les pages~7 à~11 reprennent la comonade d'une adjonction quand la base est un produit $B = \prod_{i \in I} B_i$ : $f$ est une famille $f_i : A \to B_i$, $g$ envoie $(X_i)$ sur $\prod_i g_i(X_i)$, et la page pose $\varphi_{ji} = f_j g_i$ en supposant \q{que les $f_{j}$ commutent aux produits de type $I$ (p. ex. $f_{j}$ exacts ; ou $I$ fini)} (\lot{19}{1}{7}). Les données de la comonade y deviennent (3.9) $\bar\alpha_k = \alpha_k(X_k) \circ \mathrm{pr}_k$ et $\bar\lambda_{kji} = \lambda_{kji}(X_i) \circ \mathrm{pr}_i$, avec $\lambda_{kji} : \varphi_{ki} \to \varphi_{kj}\varphi_{ji}$ venu de l'unité $\mathrm{id} \to g_j f_j$ ; la marge note que $\lambda_{iii}$ est $\lambda_i$, et le cas de deux facteurs, $g'$, $g''$ pleinement fidèles, ne laisse que $\varphi'$, $\varphi''$, $\lambda'$, $\lambda''$ (\lot{19}{1}{9}). On se donne ici des adjonctions $f_i \dashv g_i$, $f_i : A \to B_i$, et l'on suppose que $A$ a les produits indexés par $I$ ; on note $\eta_i$, $\varepsilon_i$ leurs unités et coünités, $\mathrm{pr}_i$ les projections de $\prod_i g_i(X_i)$, et $\lambda_{kji}(X) = f_k(\eta_{j, g_i X}) : f_k g_i X \to f_k g_j f_j g_i X$.

\begin{enonce}[Pages 7 et 8, la matrice de la comonade\piste{19-comonad-matrix-product-base}]
\begin{enumerate}
\item Le foncteur $f = (f_i) : A \to \prod_i B_i$ est adjoint à gauche de $g$, $g((X_i)) = \prod_i g_i(X_i)$ ; l'unité a pour composantes les $\eta_i$, et la coünité en $(X_i)$ a pour composante $j$ la flèche $\varepsilon_{j, X_j} \circ f_j(\mathrm{pr}_j)$.
\item Si $f_j$ commute aux produits indexés par $I$, la composante $j$ de $\varphi((X_i))$ est un produit des $\varphi_{ji}(X_i)$, de projections les $f_j(\mathrm{pr}_i)$.
\item Sans autre hypothèse, la coünité de $\varphi$ a pour composante $k$ la flèche $\varepsilon_{k, X_k} \circ f_k(\mathrm{pr}_k)$, et la composante $k$ de la comultiplication $\delta$ vérifie, pour tous $j$ et $i$,
\[
f_k g_j f_j(\mathrm{pr}_i) \circ f_k(\mathrm{pr}_j) \circ \delta_k = \lambda_{kji}(X_i) \circ f_k(\mathrm{pr}_i) :
\]
c'est (3.9). Si tous les $f_j$ commutent aux produits indexés par $I$, ces égalités caractérisent $\delta_k$.
\item Pour $i = j = k$, $\lambda_{iii}$ est la comultiplication de la comonade $\varphi_i = f_i g_i$.
\end{enumerate}
\end{enonce}

\begin{proof}
(1) La bijection $\operatorname{Hom}(fY, (X_i)) = \prod_i \operatorname{Hom}(f_i Y, X_i) \simeq \prod_i \operatorname{Hom}(Y, g_i X_i) \simeq \operatorname{Hom}(Y, \prod_i g_i X_i)$ est naturelle en $Y$ et en $(X_i)$, composante par composante ; l'unité est l'image de l'identité, c'est-à-dire la flèche de composantes $\eta_{i, Y}$, et la coünité l'image réciproque de l'identité, de composante $j$ la transposée de $\mathrm{pr}_j$, soit $\varepsilon_{j, X_j} \circ f_j(\mathrm{pr}_j)$.

(2) Le cône $(\mathrm{pr}_i)$ est un produit dans $A$, et $f_j$ envoie un produit sur un produit : le cône $(f_j(\mathrm{pr}_i))$ de sommet $f_j(\prod_i g_i X_i)$ est un produit de la famille $f_j g_i(X_i)$.

(3) La coünité de $\varphi$ est celle de l'adjonction, calculée en~(1). La composante $k$ de $\delta$ est $f_k(\eta_{gX})$, où $\eta$ est l'unité de~(1), dont la composante $j$ est $\eta_j$. Donc $\delta_k$ suivie de $f_k(\mathrm{pr}_j)$ puis de $f_k g_j f_j(\mathrm{pr}_i)$ vaut $f_k\bigl(g_j f_j(\mathrm{pr}_i) \circ \eta_{j, gX}\bigr) = f_k\bigl(\eta_{j, g_i X_i} \circ \mathrm{pr}_i\bigr)$, par naturalité de $\eta_j$, c'est-à-dire $\lambda_{kji}(X_i) \circ f_k(\mathrm{pr}_i)$. Si les $f_j$ commutent aux produits indexés par $I$, $f_k(\prod_j g_j f_j g X)$ est un produit de projections $f_k(\mathrm{pr}_j)$ ; $g_j$, adjoint à droite, commute aux produits, si bien que $f_k g_j f_j(\prod_i g_i X_i)$ est un produit de projections $f_k g_j f_j(\mathrm{pr}_i)$. Deux flèches qui ont les mêmes composées avec toutes ces projections sont égales.

(4) $\lambda_{iii}(X) = f_i(\eta_{i, g_i X})$ est, par définition, la comultiplication de $f_i g_i$.
\end{proof}

\begin{enonce}[Pages 7 à 9, les cas dégénérés et les deux facteurs\piste{19-comonad-matrix-product-base}]
Supposons les $g_i$ pleinement fidèles.
\begin{enumerate}
\item $\varepsilon_i : \varphi_{ii} \to \mathrm{id}_{B_i}$ est un isomorphisme ; $\lambda_{kki}$ et $\lambda_{kjj}$ sont des isomorphismes pour tous $i, j, k$.
\item Pour $i \neq j$, $\lambda_{iji} = \lambda' \circ \varepsilon_i$, où $\lambda' = \lambda_{iji} \circ \varepsilon_i^{-1} : \mathrm{id}_{B_i} \to \varphi_{ij}\varphi_{ji}$.
\item Pour deux facteurs, $B = B' \times B''$, $\varphi' = f' g''$, $\varphi'' = f'' g'$ : si $f'$ commute aux produits binaires, la composante $B'$ de $\varphi(X', X'')$ est un produit de $X'$ et de $\varphi'(X'')$, de projections $\varepsilon'_{X'} \circ f'(\mathrm{pr}')$ et $f'(\mathrm{pr}'')$ ; de même, si $f''$ y commute, la composante $B''$ est un produit de $\varphi''(X')$ et de $X''$.
\end{enumerate}
\end{enonce}

\begin{proof}
(1) Un adjoint à droite pleinement fidèle a une coünité inversible. L'identité triangulaire $\varepsilon_{k, f_k Y} \circ f_k(\eta_{k, Y}) = 1$ montre alors que $f_k(\eta_{k, Y})$ est l'inverse de $\varepsilon_{k, f_k Y}$ ; avec $Y = g_i X$, c'est $\lambda_{kki}(X)$. L'autre identité, $g_j(\varepsilon_{j, X}) \circ \eta_{j, g_j X} = 1$, montre que $\eta_{j, g_j X}$ est l'inverse de $g_j(\varepsilon_{j, X})$, et $\lambda_{kjj}(X) = f_k(\eta_{j, g_j X})$ est inversible.

(2) C'est la définition de $\lambda'$.

(3) Par~(2) de l'énoncé précédent, la composante $B'$ de $\varphi(X', X'')$ est un produit de $f'g'(X')$ et de $f'g''(X'')$ ; remplacer la première projection par sa composée avec l'isomorphisme $\varepsilon'_{X'}$ donne encore un produit, de facteurs $X'$ et $\varphi'(X'')$. De même pour $B''$.
\end{proof}

\begin{enonce}[Page 7, « ou $I$ fini »\piste{19-comonad-matrix-product-base}]
Il existe des adjonctions $f_i \dashv g_i$, avec $I$ à deux éléments, pour lesquelles la composante $j$ de $\varphi((X_i))$ n'est isomorphe au produit des $\varphi_{ji}(X_i)$ pour aucun choix d'isomorphisme.
\end{enonce}

\begin{proof}
Prenons $I = \{0, 1\}$, $A = B_0 = B_1$ la catégorie des ensembles, $f_i(Y) = \{0,1\} \times Y$ et $g_i(X) = X^{\{0,1\}}$, l'adjonction étant la curryfication, et $X_0 = X_1 = \ast$. Alors $g_i(\ast) = \ast$, $\prod_i g_i(X_i) = \ast$ et la composante $j$ de $\varphi(X)$ est $\{0,1\} \times \ast$, à deux éléments, tandis que $\prod_i \varphi_{ji}(X_i) = (\{0,1\} \times \ast)^2$ en a quatre.
\end{proof}

\begin{remarque}
L'hypothèse de la page, que les $f_j$ commutent aux produits indexés par $I$, est celle qui sert, et elle est suffisante. Des deux exemples qu'elle en donne, le premier est juste --- un foncteur exact à gauche commute aux produits finis --- mais le second ne l'est pas : un adjoint à gauche ne commute pas en général aux produits finis, et l'exemple ci-dessus en est un pour lequel la formule $\mathrm{pr}_j\,\varphi((X_i)) = \prod_i \varphi_{ji}(X_i)$ est fausse avec $I$ fini. La lecture reprend la parenthèse de la page (« ce qui est automatique si $I$ est fini, ou si les $f_j$ sont exacts ») ; il faut l'exactitude, ou directement la commutation aux produits.
\end{remarque}

\begin{remarque}
Les formules (3.9) ne demandent aucune hypothèse : elles décrivent les composées de la coünité et de la comultiplication avec les projections, et c'est seulement pour en déduire que ces données \emph{sont} une matrice, c'est-à-dire qu'elles déterminent $\varphi$ et $\delta$, qu'il faut la commutation des $f_j$ aux produits. Les cas \q{dégénérés} de la page 8, $i = j$ ou $j = k$, sont, pour des $g_i$ pleinement fidèles, ceux où $\lambda_{kji}$ est un isomorphisme ; c'est ce qui ne laisse, pour deux facteurs, que $\lambda'$ et $\lambda''$. Pour la formule $\varphi(X', X'') = (X' \times \varphi'(X''), \varphi''(X') \times X'')$, il faut encore que $f'$ et $f''$ commutent aux produits binaires, ce que la page tient de son hypothèse générale.
\end{remarque}

\begin{enonce}[Pages 14 et 15\lie{4}]
Soient $C$ une catégorie de Grothendieck, $U$ un générateur de $C$ et $R = \mathrm{End}_C(U)$. Le foncteur $\operatorname{Hom}_C(U, -) : C \to \mathbf{Mod}_R$, à valeurs dans les $R$-modules à droite, est pleinement fidèle, et son adjoint à gauche $- \otimes_R U$ est exact.
\end{enonce}

\begin{proof}
La fidélité de $\operatorname{Hom}_C(U, -)$ est la définition d'un générateur. La pleine fidélité et l'exactitude à gauche de $- \otimes_R U$ sont le théorème de Gabriel--Popescu (1964), classique, que la preuve vérifiée invoque ; l'exactitude à droite vient de ce que $- \otimes_R U$ est un adjoint à gauche.
\end{proof}

\begin{remarque}
Sous la forme de la marge, a) n'utilise que la condition~C, et non la conservativité ; la suffisance de b) est le théorème de comonadicité de Beck.
\end{remarque}

\begin{remarque}
La nécessité de b), que la page affirme par \q{il faut et il suffit}, est fausse : l'exemple ci-dessus est une adjonction comonadique dont l'adjoint à gauche ne vérifie pas la condition~C. L'erreur est sur la page, et la condition n'est pas le théorème de Beck sous sa forme précise : celui-ci caractérise les adjonctions comonadiques par une condition nécessaire et suffisante, que b) n'est pas.
\end{remarque}

\begin{remarque}
Toute double flèche $f$-scindée est une double flèche dont l'image a un noyau (lemme~\ref{lem:19-scinde}), et non l'inverse : la condition~C est plus forte que celle de Beck, suffisante sans être nécessaire.
\end{remarque}

\begin{remarque}
Pour un générateur seul, la page affirme une équivalence $C \simeq \mathbf{Mod}_R$ ; seul le plongement pleinement fidèle est acquis, et c'est la forme que la lecture retient.
\end{remarque}

\noindent Ne sont pas démontrés ici : la Conclusion des pages~4 à~6 (adjonctions vérifiant la condition~C et comonades), que la même réserve atteint ; la reconstruction de l'adjonction à partir d'une comonade ; l'énoncé de Morita (équivalence si et seulement si $U$ est petit et projectif) ; les relations de coünité et de coassociativité entre les $\lambda_{kji}$, que la page renonce à écrire, et la description de $A$ par les quadruplets $(X', X'', u', u'')$ et les triangles (3.16) (pages~9 et~10)\piste{19-comonad-matrix-product-base} ; les chaînes de $n$ adjoints successifs entre topos de préfaisceaux (pages~79 et~80)\piste{19-adjoint-chains-length-n}.
